\chapter{Geometry Profiles}\label{ch06:chapter}

\section{Why geometry, and under what discipline}\label{ch06:sec:intro}

Chapters~\ref{ch04:chapter} and~\ref{ch05:chapter} built a world of exact
identities: tagged values, canonical payloads, digests that separate
everything distinct. This chapter introduces the candidate's two
\emph{geometric} vocabularies --- the mixed-simplex Fisher--Rao baseline
and the optional projector research profile --- and, with equal care, the
disciplines that keep geometric numbers from being mistaken for exact
identities or for empirical findings.

The motivating need is real. A simplex address
(Definition~\ref{ch04:def:simplex}) is a strictly positive normalized
profile over the forty canonical ND-cells; analyses over such profiles
want a notion of \emph{closeness}: how far did a profile drift between
two ticks, how dissimilar are two declared design profiles, is this
fixture's output profile nearer to one reference than another. A
digest cannot answer these questions --- it separates everything and
grades nothing. A metric can. Information geometry supplies the
canonical candidate metric on the open simplex, with a mathematical
pedigree \cite{fisher,rao,amari,ay} that Section~\ref{ch06:sec:primer}
reviews.

But the candidate's evidence discipline requires that every geometric
capability arrive with its limits attached, and this chapter is
structured around three of them, stated here in advance:

\begin{enumerate}
\item The implemented Fisher--Rao distance is a \emph{binary64 numerical
pseudometric at its resolution}. It is disclosed as such, profiled as
such, and it is \emph{not} a proof that all distinct profiles are
separated (Section~\ref{ch06:sec:pseudometric}).
\item The projector profile is a \emph{separate, optional research
profile}, never a canonical PCBA address; its Gram-kernel overlap result
requires common rank $d > 0$ and establishes an algebraic relation
between declared artifacts, \emph{not} latent cognitive capacity
(Sections~\ref{ch06:sec:projector}--\ref{ch06:sec:gram}).
\item The CAS capacity transfer retains a \emph{proof gap} governed
solely by the formal-mathematics Paper~3 record; every engineering-side
statement about it, including everything in this chapter, is a read-only
dependency summary with no authority to narrow, close, or reinterpret
the gap (Section~\ref{ch06:sec:proofgap}).
\end{enumerate}

And over all of it stands the boundary restated in
Section~\ref{ch06:sec:boundary}: geometry profiles do not identify
learned model states with CAS states.

\section{Primer: information geometry on the open simplex}
\label{ch06:sec:primer}

\subsection{The open simplex as a statistical manifold}

Identify a simplex address $p \in \Delta_{39}^{\circ}$ with the
categorical probability distribution it defines over the forty canonical
cells: cell $k$ receives probability $p_k > 0$. The open simplex is then
the interior of the parameter space of the categorical family --- a
smooth manifold of dimension $39$, with global coordinates given by any
thirty-nine of the forty weights. Its tangent space at $p$ is the set of
zero-sum vectors,
\[
T_p \Delta_{39}^{\circ}
 = \Bigl\{\, u \in \mathbb{R}^{40} : \textstyle\sum_k u_k = 0 \,\Bigr\},
\]
reflecting that perturbations must preserve normalization.

Statistical manifolds carry a distinguished Riemannian metric, obtained
by Fisher's information quantity \cite{fisher} read as a quadratic form
on parameter perturbations; Rao's foundational observation \cite{rao}
was that this form defines a Riemannian geometry on the parameter space,
whose geodesic distance measures the intrinsic distinguishability of
nearby distributions. On the categorical family the construction is
completely explicit, and it is worth carrying it out once, because the
result is exactly the quadratic form the candidate uses.

Let $X$ be a categorical random variable over the forty cells with
distribution $p$, and perturb $p$ along a zero-sum direction $u$:
$p(t) = p + t u$, which stays in the open simplex for small $t$. The
log-likelihood of outcome $k$ is $\log p_k(t)$, its $t$-derivative at
$t = 0$ is $u_k / p_k$, and the Fisher information of the
one-parameter family at $t = 0$ is the variance of that score:
\[
I(0)
 = \sum_{k=1}^{40} p_k \Bigl( \frac{u_k}{p_k} \Bigr)^{\!2}
 - \Bigl( \sum_{k=1}^{40} p_k \frac{u_k}{p_k} \Bigr)^{\!2}
 = \sum_{k=1}^{40} \frac{u_k^2}{p_k},
\]
the cross term vanishing because $\sum_k u_k = 0$. Polarizing the
quadratic form gives the metric.

\begin{definition}[Fisher--Rao metric on the open simplex]
\label{ch06:def:frmetric}
For $p \in \Delta_{39}^{\circ}$ and
$u, v \in T_p \Delta_{39}^{\circ}$,
\[
g^{\mathrm{FR}}_p(u, v) = \sum_{k=1}^{40} \frac{u_k\, v_k}{p_k}.
\]
\end{definition}

The reasons information geometry --- rather than an arbitrary embedding
metric --- is the right setting are threefold, and each matters to how
the candidate is allowed to talk about its numbers.

\emph{Invariance.} The Fisher--Rao metric is, up to an overall scale,
the unique Riemannian metric on families of distributions invariant
under the natural statistical maps (sufficient statistics, refinement
and coarsening of outcomes); this characterization and its history are
developed in the standard treatments \cite{amari,ay}. Invariance means
the geometry measures properties of the \emph{distributions}, not of an
arbitrary coordinate choice --- exactly the property one wants if
profile comparisons are ever to be more than numerology.

\emph{Degeneracy at the boundary.} The metric coefficients $1/p_k$
diverge as $p_k \to 0$: the closed simplex is not a Riemannian manifold
under $g^{\mathrm{FR}}$. This is the geometric face of the strict
positivity rule of Chapter~\ref{ch04:chapter},
Section~\ref{ch04:sec:positivity}: the type system admits exactly the
points where the geometry is defined.

\emph{Distinguishability semantics.} Locally, the squared Fisher--Rao
length of a perturbation measures how statistically distinguishable the
perturbed distribution is from the original --- large distance means
easy to tell apart from samples, small distance means hard. This gives
distance values an operational reading \emph{as statements about
distributions}. It gives them no reading whatsoever as statements about
cognition; the distributions here are declared software profiles.

\subsection{The sphere embedding and the geodesic distance}

The categorical Fisher--Rao geometry admits a classical closed form.

\begin{proposition}[Sphere embedding]\label{ch06:prop:sphere}
The map $\varphi(p) = 2\sqrt{p} = (2\sqrt{p_1}, \ldots, 2\sqrt{p_{40}})$
is a diffeomorphism from $\Delta_{39}^{\circ}$ onto the open positive
orthant of the sphere of radius $2$ in $\mathbb{R}^{40}$, and it is an
isometry between $g^{\mathrm{FR}}$ and the metric induced on that sphere
by the Euclidean ambient metric.
\end{proposition}

\begin{proof}
If $x = \varphi(p)$, then $\sum_k x_k^2 = 4 \sum_k p_k = 4$, so the
image lies on the radius-$2$ sphere, with all coordinates positive; the
inverse $p_k = x_k^2 / 4$ is smooth on the open orthant. For the
isometry, differentiate: $dx_k = dp_k / \sqrt{p_k}$, hence
$\sum_k dx_k^2 = \sum_k dp_k^2 / p_k = g^{\mathrm{FR}}_p(dp, dp)$.
\end{proof}

\begin{corollary}[Geodesic distance]\label{ch06:cor:distance}
For $p, q \in \Delta_{39}^{\circ}$,
\[
d_{\mathrm{FR}}(p, q)
 = 2 \arccos\!\Bigl( \sum_{k=1}^{40} \sqrt{p_k\, q_k} \Bigr),
\]
the great-circle distance between $\varphi(p)$ and $\varphi(q)$ on the
radius-$2$ sphere. The argument of the arccosine --- the Bhattacharyya
coefficient of $p$ and $q$ --- lies in $(0, 1]$ on the open simplex, so
$d_{\mathrm{FR}}(p, q) \in [0, \pi)$, with $0$ attained exactly at
$p = q$.
\end{corollary}

\begin{proof}
The image of the open simplex is a geodesically convex subset of the
open positive orthant of the sphere (any two points of the open orthant
are joined by a unique minor great-circle arc lying in the orthant), so
the induced intrinsic distance is the ambient great-circle distance:
radius times the angle, and the angle between unit directions
$\varphi(p)/2$ and $\varphi(q)/2$ has cosine
$\sum_k \sqrt{p_k}\sqrt{q_k}$. By Cauchy--Schwarz this coefficient is at
most $1$ with equality iff $\sqrt{p}$ and $\sqrt{q}$ are proportional,
which under normalization forces $p = q$; positivity of all coordinates
makes the coefficient strictly positive, bounding the distance strictly
below $\pi$.
\end{proof}

\begin{remark}[Geodesic paths]\label{ch06:rem:geodesic}
The embedding also yields the geodesic itself, not only its length. For
$p \ne q$ with angle
$\alpha = \arccos\bigl(\sum_k \sqrt{p_k q_k}\bigr)$, the unit-speed
great-circle arc between $\varphi(p)$ and $\varphi(q)$ pulls back to
\[
\gamma(t)_k
 = \Bigl(
 \frac{\sin\bigl((1 - t)\,\alpha\bigr)}{\sin \alpha} \sqrt{p_k}
 + \frac{\sin(t\,\alpha)}{\sin \alpha} \sqrt{q_k}
 \Bigr)^{\!2},
 \qquad t \in [0, 1],
\]
which remains in the open simplex throughout (both coefficients are
positive for $t \in (0,1)$, and squares of positive combinations of
positive coordinates are positive; normalization follows from the arc
lying on the sphere). Interpolation between profiles, where an analysis
wants it, therefore has a canonical form --- and it is \emph{not}
the straight-line mixture $(1-t)p + tq$, which is a geodesic of a
different (flat) geometry. The candidate registers no interpolation
operation in v0.1; the formula is recorded so that any future one has
its exact target stated in advance.
\end{remark}

\begin{remark}[Neighboring distances and local behavior]
\label{ch06:rem:local}
Writing $H^2(p, q) = 1 - \sum_k \sqrt{p_k q_k}$ for the squared
Hellinger distance, the exact relation
$\cos\bigl(d_{\mathrm{FR}}(p,q)/2\bigr) = 1 - H^2(p,q)$ holds, and for
close profiles the expansion $\arccos(1 - x) = \sqrt{2x}\,(1 + O(x))$
gives
\[
d_{\mathrm{FR}}(p, q) = 2\sqrt{2}\, H(p, q)\,\bigl(1 + O(H^2)\bigr),
\qquad
d_{\mathrm{FR}}(p, p + u)^2
 = \sum_k \frac{u_k^2}{p_k} + O(\lVert u \rVert^3),
\]
the latter being the infinitesimal form one expects from
Definition~\ref{ch06:def:frmetric}: to second order, squared geodesic
distance is the chi-square-like quadratic form of the perturbation.
These classical relations \cite{amari,ay} matter operationally in
Section~\ref{ch06:sec:resolution}, where the square-root behavior of
$\arccos$ near $1$ is exactly what amplifies floating-point error for
nearby profiles.
\end{remark}

\begin{example}[A worked exact distance]\label{ch06:ex:worked}
Illustrative computation, exact mathematics only. Let $p$ be the
uniform address ($p_k = 1/40$ for all $k$) and let $q$ move mass
$0.005$ from cell \texttt{P0} to cell \texttt{P1}:
$q_1 = 0.02$, $q_2 = 0.03$, $q_k = 0.025$ otherwise. Then
\[
\sum_k \sqrt{p_k q_k}
 = 38 \cdot 0.025 + \sqrt{0.025 \cdot 0.02} + \sqrt{0.025 \cdot 0.03}
 \approx 0.95 + 0.022361 + 0.027386 = 0.999747,
\]
so $d_{\mathrm{FR}}(p, q) = 2 \arccos(0.999747) \approx 0.0450$, while
the Hellinger approximation of Remark~\ref{ch06:rem:local} gives
$2\sqrt{2}\, H \approx 2\sqrt{2} \cdot 0.015912 \approx 0.0450$ ---
agreement to the displayed precision, as expected at this scale. The
computed magnitude has, as always, no cognitive reading; it is a
distance between two declared profiles.
\end{example}

\begin{remark}[Exact-mathematics status]\label{ch06:rem:exact}
Definition~\ref{ch06:def:frmetric} through
Corollary~\ref{ch06:cor:distance} are classical mathematics
\cite{rao,amari,ay}, stated over the reals, and at that level
$d_{\mathrm{FR}}$ is a genuine metric on $\Delta_{39}^{\circ}$:
symmetric, positive on distinct pairs, and satisfying the triangle
inequality (as a geodesic distance). \emph{None of this transfers
automatically to the implementation.} The candidate's discipline is to
keep three layers separate: the exact mathematics above; the exact
Decimal data of Chapter~\ref{ch04:chapter}; and the binary64 numerical
procedure of the next section, which is where the metric properties are
partially lost and where the disclosure lives.
\end{remark}

\section{The implemented baseline: a disclosed numerical pseudometric}
\label{ch06:sec:pseudometric}

\subsection{What the implementation does}

The implemented baseline, under the versioned profile name
\texttt{cas-pcba-v8.binary64-fisher-rao-pseudometric.v1}, accepts two
exact mixed-simplex addresses and proceeds as follows: both operands are
revalidated; each exact Decimal weight is converted to binary64; the
Bhattacharyya coefficient is accumulated by compensated floating
summation of the terms $\sqrt{a_k b_k}$; the accumulated coefficient is
clamped into $[0, 1]$; the distance is computed as
$2 \arccos(\cdot)$ in binary64; results of magnitude below $10^{-15}$
are flushed to exactly zero; and the result is returned as a Decimal
constructed from the binary64 value's decimal representation. The
function object itself carries two disclosure attributes: its profile
name, and the flag \texttt{is\_exact\_metric = False}.

Every step of that pipeline is deliberate. Revalidation keeps the
garbage-in case out. The clamp exists because accumulated rounding can
push the coefficient marginally above $1$, where the arccosine is
undefined; clamping resolves the exception in the conservative
direction (distance $0$ rather than a domain error). The flush-to-zero
acknowledges that distances at the $10^{-15}$ scale are numerical noise
at binary64 resolution and refuses to report them as structure. And the
Decimal return type does \emph{not} launder the number into exactness:
it is an exact representation \emph{of the binary64 result}, tagged by
the profile that produced it.

\subsection{The disclosure, stated precisely}

\begin{definition}[Numerical pseudometric at resolution]
\label{ch06:def:pseudometric}
Call the implemented map
$\widehat{d} : \Delta \times \Delta \to \mathbb{D}$ (where $\Delta$ is
the set of valid mixed-simplex addresses and $\mathbb{D}$ the finite
decimals) a \emph{numerical pseudometric at binary64 resolution},
meaning: $\widehat{d}$ is symmetric and nonnegative by construction;
$\widehat{d}(p, p) = 0$ for every $p$; but $\widehat{d}(p, q) = 0$ does
\emph{not} imply $p = q$, and the triangle inequality is inherited from
the exact mathematics only up to the accumulated floating-point and
flush errors of the pipeline, for which no certified bound is currently
registered.
\end{definition}

\begin{finding}[Separation fails below resolution]
\label{ch06:find:separation}
There exist pairs of \emph{distinct} valid addresses --- distinct as
exact Decimal profiles, with distinct canonical identity digests ---
whose implemented distance is exactly $0$. Two mechanisms suffice:
(i)~two distinct sixty-digit Decimal weight profiles can round to the
same forty binary64 values, collapsing the computation entirely; and
(ii)~profiles distinct even after rounding can yield distances below
the $10^{-15}$ flush threshold, which the implementation reports as
zero by design. The implementation's own documentation states the
consequence: distinct profiles below binary64 resolution can have
distance zero, and the function ``is not an exact metric, proof of
identity, or projector bridge.''
\end{finding}

It must be stressed that Finding~\ref{ch06:find:separation} is a
\emph{disclosure}, not a defect report: the behavior is documented,
versioned into the profile name itself (\texttt{binary64}, \texttt{pseudometric}),
and guarded by the claim discipline below. The alternative --- exact
Decimal transcendental evaluation with certified rounding --- is a
legitimate open engineering item, and until it exists the candidate
prefers a disclosed pseudometric to an undisclosed pretension of
exactness. This is the same posture the seed specification takes:
the binary64 implementation ``is a numerical pseudometric at its
resolution --- not a proof that all distinct profiles are separated.''

\subsection{Resolution arithmetic}\label{ch06:sec:resolution}

How large are the effects the disclosure discloses? No certified error
bound is registered, and this subsection does not create one; what it
offers is a first-order \emph{heuristic} sizing, clearly labeled as
such, so that readers of reported distances can calibrate their
skepticism quantitatively rather than vaguely.

Binary64 arithmetic resolves relative differences of order
$\varepsilon \approx 2.2 \times 10^{-16}$ (the unit roundoff). Three
stages of the pipeline interact with this floor differently.

\emph{Conversion.} Each sixty-digit Decimal weight is rounded to
binary64, a relative perturbation of order $\varepsilon$ per weight.
Profiles differing only beyond the seventeenth significant digit of
every weight become \emph{identical} after conversion --- the first
mechanism of Finding~\ref{ch06:find:separation}.

\emph{Accumulation.} The Bhattacharyya coefficient is a sum of forty
terms of typical magnitude $1/40$; compensated summation keeps the
accumulated error near a few units of $\varepsilon$, so the computed
coefficient carries an absolute uncertainty on the order of
$10^{-16}$ to $10^{-15}$.

\emph{Amplification by the arccosine.} Near coincidence the coefficient
approaches $1$, where $\arccos$ has unbounded derivative: by
Remark~\ref{ch06:rem:local}, a coefficient error of magnitude $e$ near
$1$ moves the distance by roughly $2\sqrt{2e}$. Taking
$e \approx 10^{-16}$ gives a distance uncertainty of order
$2\sqrt{2 \times 10^{-16}} \approx 3 \times 10^{-8}$.

The heuristic conclusion --- and it is a heuristic, pending the
independent numerical audit of
Threat~\ref{ch06:threat:numeric} --- is that computed distances for
\emph{nearly coincident} profiles are noise-dominated below roughly the
$10^{-8}$ scale, even though the implementation only flushes below
$10^{-15}$: the flush targets exact-coincidence roundoff residue, not
the amplified noise floor near coincidence. A reported distance of,
say, $3 \times 10^{-9}$ under profile v1 should therefore be read as
``indistinguishable from zero at resolution,'' and
Rule~\ref{ch06:rule:zero}'s discipline applies to it in spirit even
though its printed value is positive. Far from coincidence the
arccosine's derivative is moderate and the same heuristic assigns
reported distances an absolute uncertainty many orders of magnitude
below their values; the discipline of
Rule~\ref{ch06:rule:positive} --- report the number, its profile, and
no exactness pretension --- covers both regimes uniformly.

\begin{example}[A binary64 collision pair]\label{ch06:ex:collision}
Illustrative construction of the first mechanism of
Finding~\ref{ch06:find:separation}. Let $p$ be the uniform address with
every weight \texttt{Decimal("0.025")}. Build $q$ from $p$ by
transferring mass $10^{-30}$ from cell \texttt{P0} to cell \texttt{P1}
at the exact-Decimal layer, keeping the exact sum at
\texttt{Decimal("1")}. Then $p$ and $q$ are distinct values with
distinct canonical identity digests --- the digest layer separates them
--- but each perturbed weight rounds to the \emph{same} binary64 number
as $0.025$, since the relative perturbation, $4 \times 10^{-29}$, is
far below the unit roundoff. The implemented pipeline computes from
identical forty-tuples of doubles, obtains a coefficient that clamps to
$1$, and reports distance exactly $0$. Both layers behaved exactly as
specified; the example is why the specification insists they never be
conflated.
\end{example}

\subsection{What this means for evidence claims}
\label{ch06:sec:evidenceclaims}

The pseudometric disclosure induces a small calculus of which claims a
Fisher--Rao number can and cannot support. The candidate states it as
rules, because the failure mode --- quietly upgrading a numerical zero
into an identity claim --- is exactly the kind of drift the audit
discipline exists to catch.

\begin{ruleenv}[Zero distance licenses nothing about identity]
\label{ch06:rule:zero}
A computed $\widehat{d}(p, q) = 0$ supports only: \emph{the implemented
pipeline could not distinguish $p$ and $q$ at binary64 resolution under
profile v1}. Identity of addresses is claimed, if at all, by canonical
digest equality (Chapter~\ref{ch04:chapter},
Rule~\ref{ch04:rule:identity}) --- the exact layer, where separation
holds by construction. Any artifact that records a zero distance and
needs an identity conclusion must record the digest comparison
alongside it.
\end{ruleenv}

\begin{ruleenv}[Positive distance carries its resolution]
\label{ch06:rule:positive}
A computed $\widehat{d}(p, q) = \delta > 0$ supports:
\emph{the profiles are distinguishable by the implemented pipeline, with
reported magnitude $\delta$ under profile v1}. Reported magnitudes are
comparable with one another under the same profile; they are not
certified approximations of the exact $d_{\mathrm{FR}}$ to any stated
error bound, because no such bound is currently registered. Downstream
text must therefore not present $\delta$ as the exact Fisher--Rao
distance, and must not chain triangle-inequality arguments through
$\widehat{d}$ as if certified.
\end{ruleenv}

\begin{ruleenv}[Profile identity travels with the number]
\label{ch06:rule:profiletravel}
Every reported distance carries the profile name
\texttt{cas-pcba-v8.binary64-fisher-rao-pseudometric.v1}; the operand
identities (their profile digests, which embed the decimal-context
digest); and, where the number enters a trace, the standard trace
bindings of Chapter~8. A distance whose provenance cannot be replayed
is not evidence of anything.
\end{ruleenv}

\begin{protocol}[Reporting template for baseline distances]
\label{ch06:proto:report}
Any artifact of this research program that reports a Fisher--Rao
baseline value records, together: (1)~the value as returned (a Decimal
transcription of the binary64 result); (2)~the profile name
\texttt{cas-pcba-v8.binary64-fisher-rao-pseudometric.v1} and its
\texttt{is\_exact\_metric = false} disclosure; (3)~the operand profile
digests; (4)~the decimal-context digest (already embedded in the
operand identities); (5)~where applicable, the trace and result
identities binding the computation to a run
(Chapter~\ref{ch05:chapter}, Protocol~\ref{ch05:proto:chain}); and
(6)~for values below the heuristic noise scale of
Section~\ref{ch06:sec:resolution}, an explicit
``indistinguishable at resolution'' qualifier in any prose summary.
Conclusions of profile \emph{identity} additionally record the digest
comparison of Rule~\ref{ch06:rule:zero}; conclusions of profile
\emph{difference} rest on digest inequality, with the distance
supplying magnitude only.
\end{protocol}

\begin{threat}[Numerical-validity residual]\label{ch06:threat:numeric}
Independent numerical audit of the baseline --- certified error bounds
for the conversion, summation, clamp, arccosine, and flush stages;
tolerance analysis; and adversarial search for worst-case separation
failures beyond the constructive mechanisms of
Finding~\ref{ch06:find:separation} --- remains open, as the threats
chapter of the seed paper records. Until closed, the pseudometric
framing is not a conservative courtesy; it is the strongest claim the
evidence admits.
\end{threat}

\section{Projector geometry: a separate optional research profile}
\label{ch06:sec:projector}

The second geometric vocabulary is deliberately quarantined from the
first. A \emph{projector profile} is an optional research artifact for
frontier analyses that want to reason about linear subspaces --- for
example, comparing declared feature subspaces across declared bases. It
is not part of the board baseline, no component port carries one except
under the explicit \textsc{projector\_profile} signal kind, and the type
itself enforces that it can never be mistaken for an address: the
discriminator field \texttt{canonical\_pcba\_address} is pinned to false
and construction rejects any other value.

\begin{definition}[Projector profile]\label{ch06:def:projector}
A \texttt{ProjectorProfile} binds: a profile identifier, basis
identifier, and source identifier (all safe references); a nonempty
square matrix of finite Decimals; a declared integer rank $r$ with
$1 \le r \le n$ (where $n$ is the dimension); a tolerance --- a finite
nonnegative Decimal capped by the versioned bound
$10^{-12}$; the fixed validation-profile name
\texttt{cas-pcba-v8.projector-validation.v1} and profile-kind
discriminator; the fixed research boundary; three \emph{recorded
validation residuals} computed at construction (symmetry residual
$\max_{i,j} \lvert M_{ij} - M_{ji} \rvert$, idempotence residual
$\max_{i,j} \lvert (M^2 - M)_{ij} \rvert$ with the square computed by
exact-sum Decimal arithmetic, and trace residual
$\lvert \trace M - r \rvert$); and two digests --- a matrix digest over
the canonical matrix payload, and a profile digest over the complete
identity payload.
\end{definition}

\begin{ruleenv}[Capped, versioned tolerance; recorded residuals]
\label{ch06:rule:tolerance}
Construction fails unless all three residuals are within the declared
tolerance, and the tolerance itself cannot exceed the versioned cap.
The residuals are not transient checks: they are \emph{recorded into
the identity payload}, so the profile digest commits to exactly how
close to an orthogonal projector of the declared rank the matrix was
found to be, under which tolerance, at validation time. Revalidation
recomputes the residuals and fails on any discrepancy with the recorded
values --- staleness and tampering are both detected --- and re-checks
the tolerance cap, so a profile cannot quietly loosen its own
acceptance criterion after the fact.
\end{ruleenv}

\begin{ruleenv}[One digest binds the whole geometry]
\label{ch06:rule:onedigest}
The profile digest is computed over a payload containing, together and
inseparably: the canonical matrix digest, the basis identifier and
source identifier, the dimension, the declared rank, the tolerance, the
three recorded residuals, the decimal-context digest, and the research
boundary. Consequently no constituent can be swapped independently ---
a re-based matrix, a re-declared rank, a retuned tolerance, or a
re-attributed source each produce a different profile identity. ``The
projector'' in any downstream artifact is always \emph{this matrix, in
this basis, at this dimension and rank, validated to this tolerance,
under this arithmetic} --- one digest, one object.
\end{ruleenv}

The design mirrors, in miniature, the board-identity argument of
Chapter~\ref{ch05:chapter} (Proposition~\ref{ch05:prop:digest}), and
for the same evidentiary reason: subspace claims are only as auditable
as the binding between the numbers and their declared interpretation
frame. The basis identifier is the frame; the digest makes frame
substitution detectable.

\begin{example}[A minimal validated projector, and a rejected one]
\label{ch06:ex:projector}
Illustrative. The $2 \times 2$ matrix
$\mathrm{diag}(\texttt{Decimal("1")}, \texttt{Decimal("0")})$ with
declared rank $1$ validates exactly: all three residuals are the
Decimal zero, within any admissible tolerance. The matrix
$\mathrm{diag}(\texttt{Decimal("0.5")}, \texttt{Decimal("0.5")})$ with
declared rank $1$ has trace residual zero but idempotence residual
$0.25$ --- twelve orders of magnitude above the tolerance cap --- and
is rejected: an ``average of two projectors'' is not a projector, and
the type refuses to launder one into subspace vocabulary. A declared
rank of $0$ or $3$ is rejected before any residual is computed; a
tolerance of $10^{-6}$ is rejected by the versioned cap, so no author
can admit a sloppy matrix by relaxing the acceptance criterion.
For validated $P = \mathrm{diag}(1, 0)$ and $Q = P$ under a shared
basis identifier, $\omega(P, Q) = 1$; for $Q' = \mathrm{diag}(0, 1)$,
$\omega(P, Q') = 0$ --- the two endpoints of
Lemma~\ref{ch06:lem:overlap}, realized at the smallest possible
dimension.
\end{example}

\section{The Gram-kernel overlap result and its limits}
\label{ch06:sec:gram}

The one registered computation over projector profiles is the
normalized overlap.

\begin{definition}[Projector overlap]\label{ch06:def:overlap}
For validated profiles $P$ and $Q$ with common dimension $n$, common
declared rank $r$, and \emph{identical declared basis identifier}, the
overlap is
\[
\omega(P, Q) = \frac{\trace(PQ)}{r},
\]
computed in exact Decimal arithmetic (exact-sum matrix products) after
revalidating both operands. The operation raises --- it does not
degrade --- on dimension mismatch, rank mismatch, or basis mismatch.
\end{definition}

The common-rank requirement is not bureaucratic: the normalization
divides by $r$, and the type system already excludes $r = 0$ (rank is
constrained to $[1, n]$ at construction), so the registered result is
available exactly when common rank $d > 0$ holds --- the condition the
seed specification states. What the number means, at the level of exact
linear algebra, is bounded by the following.

\begin{lemma}[Overlap bounds]\label{ch06:lem:overlap}
Let $P$ and $Q$ be exact orthogonal projectors on $\mathbb{R}^n$ of
common rank $d > 0$. Then
$0 \le \trace(PQ) \le d$, hence $\omega(P,Q) \in [0,1]$; moreover
$\omega(P,Q) = 1$ if and only if $\im P = \im Q$, and
$\omega(P,Q) = 0$ if and only if $\im P \perp \im Q$.
\end{lemma}

\begin{proof}
Using symmetry and idempotence,
$\trace(PQ) = \trace(PQQP) = \trace\bigl((QP)^{\top} (QP)\bigr)
 = \lVert QP \rVert_F^2 \ge 0$,
with equality iff $QP = 0$, i.e.\ $\im P \subseteq \ker Q$, i.e.\ the
ranges are orthogonal. For the upper bound, write
$\trace(PQ) = \trace(PQP)$ and note $0 \preceq PQP \preceq P$ in the
Loewner order, so $\trace(PQP) \le \trace(P) = d$; equality forces
$PQP = P$, hence for every unit $x \in \im P$,
$x^{\top} Q x = 1$, which for a projector means $x \in \im Q$; thus
$\im P \subseteq \im Q$, and equal finite rank upgrades inclusion to
equality.
\end{proof}

\begin{ruleenv}[What an overlap value may claim]
\label{ch06:rule:overlapclaims}
An overlap $\omega(P, Q)$ computed by the candidate supports exactly:
\emph{an exact-arithmetic algebraic relation between two validated
matrices that were declared to share a basis frame, each within its
recorded residuals of an orthogonal projector of the declared rank}.
Lemma~\ref{ch06:lem:overlap} then gives the value its geometric
reading --- subspace alignment --- \emph{up to} the residual tolerances,
whose propagation into an error bar on $\omega$ is an open numerical
item. What the value does \emph{not} support, under any magnitude: that
either subspace corresponds to a capability, feature, concept, or
state of any learned model or cognitive system. The overlap is a
statement about declared artifacts, and the registered result is
explicitly recorded as \emph{not} establishing latent cognitive
capacity.
\end{ruleenv}

The last sentence deserves its own paragraph, because it is where a
plausible-sounding inference chain must be cut. The tempting chain
runs: a projector was derived from model activations (elsewhere, under
some analysis); two such projectors overlap strongly; therefore the
model has a shared latent structure; therefore a cognitive capacity is
present. Every link after the first imports assumptions the candidate
does not admit: that the derivation basis is meaningful, that overlap
of declared artifacts measures overlap of whatever produced them, and
that subspace structure is capacity. The projector profile's fields are
designed to keep the first link honest (source and basis identifiers,
recorded residuals) precisely so that the later links can be seen to be
\emph{missing} rather than implied. The Gram-kernel result is
mathematics about matrices; capacity claims are empirical claims about
systems; the candidate registers the former and quarantines the latter.

\section{The capacity-transfer proof gap}\label{ch06:sec:proofgap}

Beyond the overlap result lies a stronger statement that the research
program cares about: a transfer principle by which structure validated
in one representation would carry CAS-semantic weight in another.
This monograph's obligation is to record its status exactly.

\begin{hypothesis}[CAS capacity transfer --- status: open, not licensed]
\label{ch06:hyp:transfer}
The proposed transfer from projector-level structure to CAS capacity
claims is a hypothesis. It retains a \emph{proof gap}: no admitted
derivation closes the distance between algebraic subspace relations
(Lemma~\ref{ch06:lem:overlap} and its registered computational form)
and any statement about cognitive or model capacity.
\end{hypothesis}

\begin{ruleenv}[Governance of the gap]\label{ch06:rule:papergov}
The proof gap is governed \emph{solely} by the formal-mathematics
Paper~3 record, at that record's own commit. Engineering records ---
the implementation, its tests, the seed research paper, and this
monograph --- are \emph{read-only dependency summaries} with respect to
that record: they may restate its status, as this section does, and
they must track it; they have no authority to close, narrow, extend, or
reinterpret it. In particular, no accumulation of green software tests,
no overlap computation of any magnitude, and no prose in an engineering
artifact changes the status of Hypothesis~\ref{ch06:hyp:transfer}. If
the Paper~3 record changes, engineering summaries update to match ---
never the reverse.
\end{ruleenv}

This governance rule is an instance of the candidate's general
authority discipline (the same one that keeps locally constructed
calibration records out of Evidence Packets and self-hashed receipts out
of licensing): the artifact that \emph{could} benefit from a claim is
never the artifact that adjudicates it. A geometry module that could
declare its own mathematics sufficient would be exactly the kind of
self-certifying loop the architecture exists to prevent.

\section{Two arithmetic regimes, and why they differ}
\label{ch06:sec:regimes}

A reader tracking the candidate's arithmetic discipline
(Chapter~\ref{ch04:chapter}, Section~\ref{ch04:sec:decimal}) will have
noticed an asymmetry in this chapter that deserves to be made explicit,
because it looks like an inconsistency and is in fact a design rule.

The \emph{projector} path is exact end to end. Residual computation
multiplies matrices with exact aligned-integer Decimal summation;
overlap traces are exact sums of exact products; the tolerance
comparison is exact Decimal comparison. Nothing in the projector path
introduces a rounding step beyond the fixed sixty-digit context's
governance of individual multiplications, and the recorded residuals
are reproducible to the digit on any conforming interpreter.

The \emph{simplex distance} path is not. It leaves the exact world at
the conversion step and returns only a transcription of a binary64
result, with the pseudometric disclosure attached.

The difference is not carelessness on one side; it is the difference
between polynomial and transcendental computation. Residuals, traces,
and overlaps are polynomial in the matrix entries: finitely many exact
multiplications and exact additions, closed under Decimal arithmetic.
The Fisher--Rao distance requires square roots and an arccosine ---
transcendental functions whose exact values are not decimals at all, so
\emph{some} approximation contract is mandatory. The candidate's v0.1
choice is the weakest honest one: use the platform's binary64
transcendentals, and disclose the resulting object as a pseudometric at
that resolution rather than dress it in exactness it does not have. A
future profile could tighten the contract --- correctly rounded
sixty-digit square roots exist in the Decimal context, and
arbitrary-precision arccosine with certified error is implementable ---
but that would be a \emph{new versioned profile} with its own name and
its own disclosure, not a quiet upgrade of v1
(Rule~\ref{ch06:rule:profiletravel} makes the version part of every
reported number precisely so that such an upgrade would be visible in
the evidence).

The same asymmetry explains an operational rule that might otherwise
puzzle: digest-bearing identity payloads embed \emph{weights}, never
\emph{distances}. A simplex address's identity contains its forty exact
Decimals; a projector's identity contains its exact matrix and exact
residuals; no identity anywhere contains a Fisher--Rao value. Exact
objects are identity material; approximate objects are reportable
results. Keeping the two segregated is what lets the candidate promise
that identities are deterministic across conforming platforms while
transcendental libraries vary.

\section{Geometry values in the wiring and the trace}
\label{ch06:sec:wiring}

Both geometry types are first-class citizens of the board vocabulary,
and their integration points are worth enumerating because each is also
a containment point.

\emph{Signal kinds and units.} A net can carry a simplex address only
under the \textsc{mixed\_simplex\_address} signal kind with unit
\textsc{probability}, and a projector profile only under
\textsc{projector\_profile} with unit \textsc{none}
(Chapter~\ref{ch04:chapter}, Table~\ref{ch04:tab:sigunit}). The closed
admissibility table means geometry values cannot ride generic scalar
wiring, and scalars cannot pose as geometry.

\emph{Port validation on canonical bytes.} When a geometry value
crosses a port, validation happens against its serialized canonical
identity: a mixed-simplex port requires a payload of mixed-simplex
type whose embedded identity revalidates (canonical cell labels, forty
tagged positive weights, exact sum one, matching decimal-context
digest); a projector port likewise requires a full projector identity
whose recorded residuals, tolerance cap, rank bounds, and boundary
block all re-check, and whose matrix-digest linkage is intact. A
tampered geometry payload does not reach a component; it fails the
port.

\emph{Measurements.} A geometry value observed at a test point enters
the trace as a measurement record whose value text is the canonical
identity payload --- so the full forty weights, or the full projector
identity, are in the evidence verbatim, under the value digest. Any
distance later computed \emph{from} traced addresses is a derived
result governed by Protocol~\ref{ch06:proto:report}; the trace holds
the exact operands, so the derivation is replayable by any auditor
with a conforming implementation.

\emph{The adapter boundary.} Neither geometry type is
finite-adapter-eligible, and the TKS adapter rejects both by type:
continuous profiles are exactly the objects the finite-arrow boundary
exists to keep out (Chapter~\ref{ch04:chapter},
Rule~\ref{ch04:rule:nottks}). The geometry chapter therefore ends
where the adapter chapter begins: with two vocabularies that touch at
zero points, by construction.

\section{Separations between the two profiles}\label{ch06:sec:separation}

It is worth closing the technical development with the separations the
two geometry profiles maintain, since their coexistence in one chapter
could otherwise suggest a unity the candidate does not claim.

\begin{invariant}[Profile separations]\label{ch06:inv:separation}
(i) A mixed-simplex address is never a projector profile, and neither
is ever coerced to the other; they occupy distinct signal kinds
(\textsc{mixed\_simplex\_address} versus \textsc{projector\_profile})
with distinct admissible units, so no net can deliver one where the
other is expected.
(ii) The Fisher--Rao baseline consumes addresses only; the overlap
consumes projector profiles only; each raises on the other's type. The
baseline's own documentation disclaims being a ``projector bridge.''
(iii) The baseline is part of the candidate's named geometry
(\texttt{cas-pcba-v8.mixed-simplex.v1} addresses;
\texttt{cas-pcba-v8.binary64-fisher-rao-pseudometric.v1} distance);
the projector profile is optional research apparatus
(\texttt{cas-pcba-v8.projector-validation.v1}) that no baseline
artifact depends on.
(iv) Both profiles bind the decimal-context digest and the fixed
research boundary into their identities; neither can be constructed in
an ``accepted'' or ``licensed'' state.
\end{invariant}

\begin{exercise}[Boundary behavior of the distance]
\label{ch06:ex:boundaryexercise}
Show that on the \emph{open} simplex the supremum $\pi$ of
$d_{\mathrm{FR}}$ is not attained: for any two valid addresses the
Bhattacharyya coefficient is strictly positive, while sequences of
addresses concentrating on disjoint cells drive it to $0$. Conclude
that ``maximally distant profiles'' is not a well-posed superlative in
the candidate's vocabulary, and relate this to the exclusion of one-hot
addresses (Chapter~\ref{ch04:chapter}, Rule~\ref{ch04:rule:notonehot}).
\end{exercise}

\begin{exercise}[Geodesic midpoint]\label{ch06:ex:midpointexercise}
Using Remark~\ref{ch06:rem:geodesic}, derive the closed form of the
geodesic midpoint of two addresses $p, q$ and verify it lies in the
open simplex; then explain why the arithmetic mixture $(p + q)/2$
generally differs from it, and compute both for the pair of
Example~\ref{ch06:ex:worked}.
\end{exercise}

\begin{remark}[Why the quarantine is architecture, not caution]
The two vocabularies answer different questions with different
mathematics: distances between normalized profiles (a Riemannian
question on a statistical manifold) versus alignment between declared
subspaces (a linear-algebra question in a declared frame). Merging them
--- for example, defining an ad hoc ``distance'' between an address and
a projector --- would require exactly the identification of profile
space with representation space that the boundary below prohibits.
The type-level quarantine makes the prohibited identification not
merely unclaimed but unrepresentable in the candidate's vocabulary,
which is the strongest form of claim discipline software can provide.
\end{remark}

\section{Claim boundary}\label{ch06:sec:boundary}

\begin{boundary}\label{ch06:bnd:main}
This chapter presents mathematics and software contracts of the CAS
PCBA~8 v0.1 research candidate, at Ring~4, \textsc{not licensed}. The
exact Fisher--Rao results
(Definition~\ref{ch06:def:frmetric}--Corollary~\ref{ch06:cor:distance})
are classical statements about probability distributions
\cite{fisher,rao,amari,ay}; their subjects here are declared software
profiles, not measured states of any system. The implemented distance
is a disclosed binary64 numerical pseudometric at its resolution and is
not a proof that all distinct profiles are separated
(Finding~\ref{ch06:find:separation}); every claim built on its values
is bounded by Rules~\ref{ch06:rule:zero}--\ref{ch06:rule:profiletravel}.
The projector Gram-kernel result requires common rank $d > 0$ and
establishes algebraic relations between declared, validated artifacts
only; it does not establish latent cognitive capacity
(Rule~\ref{ch06:rule:overlapclaims}). The CAS capacity transfer retains
a proof gap governed solely by the formal-mathematics Paper~3 record,
with respect to which all engineering records, including this chapter,
are read-only dependency summaries
(Hypothesis~\ref{ch06:hyp:transfer}, Rule~\ref{ch06:rule:papergov}).
Geometry profiles do not identify learned model states with CAS states;
no physical, causal, safety, clinical, regulatory, or product claim is
made or licensed by anything in this chapter.
\end{boundary}
